% !TEX program = pdflatex
%
% Автореферат магистерской диссертации Р. Н. Майоршина.
% Оформлен тем же классом pmk-graduate в режиме [master,summary], что и
% автореферат Антонова: класс грузит A5 twoside и сам рисует титульный лист
% (\maketitlepage). Метаданные заданы здесь же (отдельного metadata у этой
% работы нет). Содержание перенесено из исходного autoreferat_mayorshin
% (профиль «Системный анализ», петрофизическая интерпретация данных ГИС).
%
% Сборка содержания:        pdflatex autoreferat_mayorshin
% Раскладка «книжкой» (A4): pdflatex autoreferat_mayorshin_booklet
%
% В summary-режиме класс делает \section эквивалентом \section* (без номера),
% поэтому заголовки пишутся как \section{...} БЕЗ звёздочки.
%
\documentclass[master,summary]{pmk-graduate}

\usepackage{amsmath}
\usepackage{amssymb}
\setlist[enumerate]{label=\arabic*.}

% --- Метаданные работы ------------------------------------------------------
\program{msa}% 01.04.02 Прикладная математика и информатика, профиль «Системный анализ»

\title{%
Разработка метода автоматической настройки данных ГИС в рамках
петрофизической интерпретационной модели%
\english
Development of an Automatic Tuning Method for Well-Logging (LAS) Data within a
Petrophysical Interpretation Model%
}

\student{%
Майоршин Руслан Новрузович%
\brief Р.\,Н.\,Майоршин%
\group М-22%
\english Mayorshin Ruslan Novruzovich%
}

\superviser{%
Багрова Инна Александровна%
\brief И.\,А.\,Багрова%
\position %
\degree кандидат физико-математических наук%
}

\year{2025-2026}

\goal{%
Разработка метода и программного обеспечения для автоматизированной подготовки
(настройки) данных геофизических исследований скважин (ГИС) из файлов LAS к
требованиям выбранной петрофизической интерпретационной модели.%
}

\begin{document}
\maketitlepage

\section{Актуальность}

Диссертационная работа посвящена разработке метода автоматической подготовки и настройки данных геофизических исследований скважин, представленных в формате LAS, к требованиям выбранной петрофизической интерпретационной модели.

В производственной практике данные ГИС являются одним из основных источников информации о литологическом составе и свойствах разреза. Однако исходные LAS-файлы редко бывают непосредственно пригодны для автоматической интерпретации. Кривые могут иметь разные мнемоники, различные числовые диапазоны, неодинаковые единицы измерения, пропуски, локальные выбросы и участки с низким качеством записи. Даже если файл корректно читается программно, его значения не всегда согласованы с параметрами интерпретационной модели.

Ручная подготовка таких данных требует времени и зависит от опыта интерпретатора. Специалисту приходится проверять шкалы, удалять аномалии, выбирать кросс-плоты, оценивать положение облаков точек относительно модельных узлов и вручную подбирать параметры. Это снижает воспроизводимость результата и затрудняет повторную проверку интерпретации.

Поэтому актуальна разработка метода, который не просто нормализует данные статистически, а приводит кривые к пространству выбранной петрофизической модели. Такой подход позволяет сопоставлять исходные признаки с модельными параметрами, рассчитывать производные индексы, оценивать доли литотипов по глубине и сохранять результат в виде расширенного LAS-файла.

\section{Цель работы}

\goalname

\section{Задачи}

Для достижения поставленной цели в работе решались следующие задачи:
\begin{enumerate}
    \item Выполнить аналитический обзор и обосновать выбор метода автоматической настройки данных ГИС, включая его теоретические основы и область применимости.
    \item Сформулировать требования к входным данным и целевым результатам; описать общие принципы унификации, сопоставления и согласования кривых в рамках интерпретационной модели.
    \item Разработать и реализовать программный прототип, обеспечивающий загрузку исходных данных, базовую автоматическую настройку и визуализацию промежуточных и итоговых результатов.
    \item Провести апробацию подхода на репрезентативной выборке; оценить корректность и устойчивость метода по согласованным показателям и сопоставить с существующей практикой.
    \item Подготовить выводы и рекомендации по применению и дальнейшему развитию решения, включая ограничения и направления доработки.
\end{enumerate}

\section{Используемая модель и постановка задачи}

В работе используется петрофизическая интерпретационная модель, включающая восемь литотипов: песчаник, алевролит, аргиллит, мергель, известняк, доломит, ангидрит и гипс. Для каждого литотипа заданы характерные значения геохимических и петрофизических признаков: калия, тория, урана, кальция, магния, кремния, серы, плотности, гамма-каротажа, глинистости и ряда дополнительных параметров.

Модель рассматривается как система узлов в пространстве признаков. Для каждого литотипа задаётся вектор модельных значений
\[
    m_j=(M_{j,1},M_{j,2},\ldots,M_{j,F}),
\]
где $j$ --- номер литотипа, $F$ --- число признаков. Задача метода состоит в том, чтобы преобразовать исходные кривые LAS-файла к этому пространству и затем оценить доли литотипов по глубине.

На каждой глубине требуется найти вектор долей
\[
    v(d)=(v_1(d),\ldots,v_J(d)), \qquad v_j(d)\ge 0,\quad \sum_j v_j(d)=1.
\]
Эти ограничения отражают физический смысл результата: доли литотипов не могут быть отрицательными и должны суммироваться в 100\%.

\section{Метод автоматической настройки}

Разработанный метод представляет собой последовательность обработки LAS-файла. Сначала выполняется чтение файла, извлечение глубинной сетки, кривых, единиц измерения и значения отсутствующих данных. Затем мнемоники кривых сопоставляются с внутренними обозначениями признаков модели. Например, калиевый канал может быть представлен как \texttt{POTA}, \texttt{K} или \texttt{POTASSIUM}, а кальциевый --- как \texttt{CAC}, \texttt{CA} или \texttt{CALCIUM}.

Далее выполняется обработка пропусков и локальных выбросов. Для выявления аномальных точек используется робастный локальный критерий на основе медианы и медианного абсолютного отклонения. Это позволяет исключить одиночные выбросы из подбора параметров калибровки, не разрушая исходный массив данных.

В базовом режиме калибровка выполняется по устойчивым квантилям. Нижний и верхний квантили реальной кривой сопоставляются с минимальным и максимальным значениями соответствующего признака в модели. Для кривой $x_f(d)$ строится линейное преобразование
\[
    y_f(d)=a_f+b_f x_f(d),
\]
где коэффициенты выбираются так, чтобы рабочий диапазон исходной кривой совпал с диапазоном модельного параметра.

При наличии контрольных литотипных профилей применяется объектная калибровка. По контрольным долям и таблице модели строится ожидаемая кривая признака, после чего подбирается линейное или полиномиальное преобразование исходного канала. Такой режим используется не как обязательный вход метода, а как способ адаптации модели к конкретной скважине или месторождению.

После калибровки рассчитываются производные индексы:
\[
    CI=\frac{Ca}{Ca+Si+Mg}, \qquad
    DI=\frac{Mg}{Ca+Mg}, \qquad
    SI=\frac{Si}{Ca+Si+Mg}.
\]
Они помогают выделять карбонатную, доломитизированную и силикатную составляющие. Для спектральных признаков дополнительно используются отношения $Th/K$, $Th/U$ и вычисленный гамма-каротаж.

Оценка долей литотипов формулируется как регуляризованная задача неотрицательной смесевой аппроксимации:
\[
    v^*(d)=\arg\min_{v\ge0,\;\mathbf 1^Tv=1}
    \|W(M^Tv-y(d))\|_2^2+\lambda\|v-v(d-1)\|_2^2.
\]
Первый член отвечает за согласование рассчитанной смеси с наблюдаемыми признаками, второй --- за устойчивость профиля по глубине. После решения выполняются сглаживание, отсечение малых долей и повторная нормировка.

\section{Программная реализация}

Метод реализован в виде веб-системы LAS Viz. Система включает клиентскую часть, сервер API, вычислительный модуль, базу данных, очередь фоновых заданий и объектное хранилище. Клиентская часть отвечает за загрузку файлов, просмотр кривых, кросс-плоты, настройку параметров и отображение результатов. Сервер API управляет пользователями, файлами, версиями и заданиями. Вычислительный модуль выполняет расчёты на серверной стороне.

Для реализации использованы React, TypeScript, Plotly, NestJS, PostgreSQL, Redis, BullMQ и S3-совместимое хранилище. Вычислительное ядро вынесено в отдельный пакет, что обеспечивает повторное использование алгоритма и воспроизводимость расчётов.

Система поддерживает загрузку LAS-файлов, хранение исходной версии, запуск интерпретации, отображение калиброванных кривых и производных индексов, просмотр долей литотипов, хранение нескольких версий интерпретации и экспорт результата в LAS-файл. Исходный файл не перезаписывается: каждый запуск создаёт новую версию.

\section{Апробация и результаты}

Апробация выполнена на данных работающих нефтяных скважин. Для проверки использовались два LAS-файла, содержащие спектральные и геохимические кривые. Контрольные литотипные профили, полученные применяемой на практике методикой, использовались для оценки качества после расчёта.

В базовом режиме контрольные профили не использовались в качестве входных данных. Метод сравнивался с простейшим способом сопоставления каждой глубины с ближайшим узлом модели. Для первого файла средняя корреляция выросла с 0.3504 до 0.5240, а RMSE снизился с 25.03 до 16.83 процентных пункта. Для второго файла средняя корреляция выросла с 0.2893 до 0.5119, а RMSE снизился с 25.62 до 18.70 процентных пункта.

В полном режиме с объектной адаптацией RMSE снизился до 11.25 и 9.47 процентных пункта для двух файлов соответственно. Абляционный анализ показал, что основной вклад дают калибровка признаков, неотрицательная смесевая оценка, регуляризация по глубине, сглаживание и настройка узлов модели.

Анализ каналов показал, что наиболее устойчиво калибруются \texttt{POTA}, \texttt{THOR}, \texttt{CAC}, \texttt{SII} и \texttt{SCAP}. Магниевый канал оказался менее информативным в использованных данных. Кроме того, слабым местом остаётся разделение гипса и ангидрита: для более уверенной дифференциации этих литотипов необходим водородный признак.

\section{Структура работы}

Работа включает введение, четыре главы, заключение и список литературы.

В первой главе рассмотрены данные ГИС, формат LAS, петрофизическая интерпретационная модель, существующие подходы и постановка задачи. Во второй главе описан разработанный метод автоматической настройки. В третьей главе приведена архитектура программной системы и пользовательский сценарий. В четвёртой главе выполнена апробация метода и анализ результатов. В заключении сформулированы основные выводы и направления дальнейшего развития.

\section{Основные результаты}

В результате выполнения работы:
\begin{enumerate}
    \item Обоснована необходимость автоматической настройки данных ГИС перед применением петрофизической модели.
    \item Разработан метод, объединяющий робастную подготовку кривых, модельно-ориентированную калибровку, расчёт индексов и регуляризованную оценку долей литотипов.
    \item Реализована веб-система, обеспечивающая загрузку LAS, визуализацию, запуск расчёта, хранение версий и экспорт результата.
    \item Проведена апробация на данных работающих нефтяных скважин.
    \item Показано улучшение качества относительно базового способа сопоставления с ближайшим узлом модели.
    \item Выявлены ограничения метода и определены направления дальнейшей доработки.
\end{enumerate}

\section{Публикации автора по теме ВКР}

Майоршин Р.\,Н. Автоматическая настройка данных геофизических исследований скважин в рамках петрофизической интерпретационной модели // Студенческая конференция факультета ПМиК. Сборник трудов.~--- Тверь: ТвГУ, 2026.~--- С.~126--135.

\end{document}
